\documentclass[journal=tches,submission]{iacrtrans}
%%%% NOTES:
% - Change "submission" to "final" for camera-ready version
% - Add "spthm" for LNCS-like theorems if needed

%%%% 2. PACKAGES %%%%
\usepackage{amsmath,amssymb,amsthm}
\usepackage{graphicx}
\usepackage{url}
\usepackage{xcolor}

\usepackage{algorithm}
\usepackage{algorithmicx, algpseudocode}
\usepackage{tikz}
\usetikzlibrary{arrows.meta}
\usepackage{pgfplots}
\pgfplotsset{compat=1.18}

\renewcommand{\algorithmicrequire}{\textbf{Input:}}
\renewcommand{\algorithmicensure}{\textbf{Output:}}

% Custom commands for common notation
\newcommand{\F}{\ensuremath{\mathbb{F}} }
\newcommand{\Fq}{\ensuremath{\F_q} }
\newcommand{\Fe}[1]{\ensuremath{\F_{q^{#1}}} } % Field extension
\newcommand{\G}{\ensuremath{\mathbb{G}} }
\newcommand{\gComm}[1]{\color{gray}\Comment{#1}\color{black}}
\newcommand{\gText}[1]{\color{gray}#1\color{black}}

\newcommand{\IOP}[4]{\vspace{0.5cm}\noindent\fbox{\parbox{\dimexpr\linewidth-2\fboxsep-2\fboxrule\relax}{
\begin{center}
    \begin{tikzpicture}[node distance=11cm, every node/.style={font=\footnotesize},>-{>[length=9pt, width=10pt,line width=2pt]}]
    \node (L)[draw,font=\large,rectangle] {#2};\node (R)[draw,font=\large,rectangle, right of=L] {#3};
    #4
    \end{tikzpicture}
    #1
\end{center}
}}
}

\newcommand{\TikzArrow}[4]{
    \draw[gray,>-{>[length=7pt, width=8pt]},line width=1.5pt] ([yshift=-{#4}cm]#1.east) -- ([yshift=-{#4}cm]#2.west) node[midway, above, black] {#3};
}

\newcommand{\LDoes}[2]{\node[draw,rectangle,above=1pt,right,align=left] at ([yshift=-{#2}cm]L.west) {#1};}
\newcommand{\RDoes}[2]{\node[draw,rectangle,above=1pt,left,align=right] at ([yshift=-{#2}cm]R.east) {#1};}
\newcommand{\RtoL}[2]{
    \usetikzlibrary {arrows.meta}
    \draw[gray,>-{>},line width=1.5pt] ([yshift=-{#2}cm]R.east) -- ([yshift=-{#2}cm]L.west) node[midway, above, black] {#1};
}
\newcommand{\LtoR}[2]{
    \usetikzlibrary {arrows.meta}
    \draw[gray,>-{>[length=9pt, width=10pt,line width=2pt]},line width=1.5pt] ([yshift=-{#2}cm]L.west) -- ([yshift=-{#2}cm]R.east) node[midway, above, black] {#1};
}

%%%% 3. AUTHOR, INSTITUTE %%%%
\author{Shramee Srivastav\inst{1}}
\institute{
  MIST.cash, Singapore, Singapore, \email{shramee@proton.me}
}
%%%% NOTES:
% - ORCID can be added in camera-ready version if needed
% - City name is required for indexation

%%%% 4. TITLE %%%%
\title{Pairing Proofs and Polynomial Ring Toolkit}
%%%% NOTES:
% - If title is too long, provide running title: \title[Short]{Long}

\emergencystretch=1em

\begin{document}

\maketitle

%%%% 5. KEYWORDS %%%%
\keywords{Elliptic curve pairings \and Zero-knowledge proofs \and Polynomial identity testing \and Blockchain verification \and Recursive proofs.}

%%%% 6. ABSTRACT %%%%
\begin{abstract}
Elliptic curve pairings are fundamental to modern zero-knowledge proof systems, including Groth16, PLONK, and KZG polynomial commitments. In resource-constrained execution environments (Ethereum VM, BitVM, blockchain smart contracts) where computation is metered by gas fees, and in arithmetized computation systems (R1CS, Plonkish, Cairo AIR) where each computational step imposes significant prover overhead, the cost of pairing computation remains a practical bottleneck.

We present a public-coin interactive proof system that consolidates the Miller loop iterations and final exponentiation into a single polynomial relation. The relation can then be verified as polynomial identities greatly reducing the pairing costs.

We present a full implementation in gnark~\cite{gnark}---the industry-standard framework for arithmetized circuit development---as a reusable polynomial ring toolkit in the \texttt{emulated} package. Benchmarked on a Groth16 verification simulation (2 fixed-point pairings + 1 full pairing + 1 $\mathbb{F}_{p^{12}}$ multiplication) on BN254, our gnark implementation achieves a \textbf{39\% reduction in SCS constraints} (1,890,771~$\to$~1,158,194), a \textbf{60\% reduction in commitments} (739,198~$\to$~295,946), and a $\sim$\textbf{42\%} reduction in compile-time memory and up to \textbf{52\%} in solver memory. These improvements hold across both SCS and R1CS constraint systems and enable practical pairing verification in recursive proof systems and resource-constrained environments.
\end{abstract}

%%%% 7. PAPER CONTENT %%%%
\section{Introduction}

Blockchain technology continues to mature with improving scalability and infrastructure. However, privacy remains a significant barrier to entry for many users and applications. Zero-knowledge proofs (ZKPs) have emerged as a powerful solution to this challenge.

\subsection{Motivation}


Elliptic curve pairings represent fundamental building blocks for numerous cryptographic systems including Groth16\cite{Gro16}, PLONK\cite{plnk19}, BLS signature schemes\cite{bls05}, and KZG polynomial commitment schemes\cite{kzg10}.

However, \textbf{pairings are computationally intensive}. This creates critical bottlenecks in two important deployment contexts:

\begin{enumerate}
\item \textbf{Resource-constrained virtual machines} (Ethereum VM, BitVM, ZKVMs): Computation is metered — each operation incurs gas costs. A single pairing verification requiring thousands of field operations becomes economically prohibitive especially during peak network congestion.

\item \textbf{Arithmatized computation systems} (R1CS~\cite{r1cs13}, Plonkish~\cite{plnk19}, Cairo AIR~\cite{cairo21}): Each operation must be represented algebraically, adding significant prover overhead beyond native execution costs.
\end{enumerate}

In such contexts, we can make pairing verification practical by using an efficient proof system to verify the correctness of the pairing computation, rather than executing all operations natively within the expensive environment.

\subsection{Our Contribution}

We present a public-coin interactive proof system for verifying pairing computations that consolidates verification into unified polynomial relationships under the Schwartz-Zippel lemma. By treating complete Miller loop iterations as single polynomial relationships rather than sequences of individual field operations\cite{FXFM23}, our approach achieves:

\begin{itemize}
\item \textbf{45\% reduction} in Miller loop field operations
\item \textbf{77\% reduction} in commitment overhead (Fiat-Shamir hashing)
\item \textbf{37\% overall speedup} for complete pairing verification
\end{itemize}

Our key insight is that instead of verifying individual field operations separately, we can verify entire Miller loop iterations as single polynomial identities. This consolidation dramatically reduces both the number of modular operations required and the communication overhead from intermediate commitments.

\section{Background}

\begin{algorithm}
\caption{Optimal Ate pairing over Barreto-Naehrig curves~\cite{hsp10}}
\begin{algorithmic}[1]
\Require{$P \in \G_1, Q \in \G_2$}
\Ensure{$a_{opt}(Q, P)$}
\State{Write $s = 6t + 2 = \sum_{i=0}^{L-1} s_i 2^i$ where $s_i \in \{-1, 0, 1\}$ and $s_L = 1$}
\State{$T \leftarrow Q, f \leftarrow 1$}
\For{$i = L - 2$ to $0$}
    \State{$f \leftarrow f^2 \cdot l_{T,T}(P)$}
    \State{$T \leftarrow [2]T$}
    \If{$s_i = 1$}
        \State{$f \leftarrow f \cdot l_{T,Q}(P)$}
        \State{$T \leftarrow T + Q$}
    \EndIf
    \If{$s_i = -1$}
        \State{$f \leftarrow f \cdot l_{T,-Q}(P)$}
        \State{$T \leftarrow T - Q$}
    \EndIf
\EndFor
\State{$Q_1 \leftarrow \pi_p(Q), Q_2 \leftarrow \pi_p^2(Q)$}
\State{$f \leftarrow f \cdot l_{T,Q_1}(P), T \leftarrow T + Q_1$}
\State{$f \leftarrow f \cdot l_{T,-Q_2}(P), T \leftarrow T - Q_2$}
\State{$f \leftarrow f^{p^{12} - 1 / r}$}
\State{\Return $f$}
\end{algorithmic}
\end{algorithm}

\subsection{Pairing Computation}

An elliptic curve pairing is a non-degenerate bilinear map $e: \G_1 \times \G_2 \rightarrow \G_T$, where $\G_1$ and $\G_2$ are groups of points on elliptic curves over finite fields, and $\G_T$ is a multiplicative group in an extension field. For practical cryptographic applications, we require that this map be efficiently computable. We will focus on the optimal ate pairing~\cite{hsp10} over Barreto-Naehrig (BN) curves, which is widely used due to its efficiency. We will specifically use the Ethereum's BN254 curve, defined over a 254-bit prime field $\Fq$, with groups $\G_1, \G_2$ of order $r$ and target group $\G_T$ contained in the 12th extension field $\Fe{12}$.

The computation of an optimal ate pairing---the focus of our implementation---consists of two primary phases:

\begin{enumerate}
    \item \textbf{Miller loop}: The core iterative computation that evaluates line functions at pairs of points from $\G_1$ and $\G_2$, accumulating intermediate results to produce a value in $\Fe{12}$.
    
    \item \textbf{Final Exponentiation}: This step raises the Miller loop output to the power $(q^{12}-1)/r$, where $r$ is the order of the groups, ensuring the result lies in the correct subgroup $\G_T$.
\end{enumerate}

While both phases present optimization opportunities, the constraint-based verification of these computations introduces unique challenges and possibilities that differ significantly from native implementations.

\subsection{Schwartz-Zippel lemma} \label{sec:bg:sz}

The Schwartz-Zippel lemma (or more precisely, the Demillo-Lipton-Schwartz-Zippel lemma~\cite{DemLip78,Sch79,Zip79,Sch80}) provides a probabilistic method for testing polynomial identities. The lemma states that for a non-zero polynomial $P(x_1, x_2, \ldots, x_n)$ of total degree $d$ over a finite field $\F$, when evaluated at uniformly random field elements, the probability that $P$ evaluates to zero is at most $d/|\F|$.

This probabilistic bound becomes negligible for cryptographically sized fields where $d \ll |\F|$, making it practical to verify polynomial identities by evaluation instead of coefficient-wise operations.

\subsection{Constraint-Based Computation Models}

\subsubsection{Computational Model}

Our optimization targets constraint-based computation models where programs are represented as systems of algebraic constraints over finite fields. We formalize this model as follows:

\textbf{Definition (Constraint System):} A constraint system over a finite field $\Fq$ consists of:
\begin{itemize}
\item A set of variables $\mathbf{v} = (v_1, \ldots, v_m) \in \Fq^m$
\item A partition of variables into:
\begin{itemize}
\item \textit{Public inputs} $\mathbf{x} \in \Fq^n$ (known to verifier)
\item \textit{Private witnesses} $\mathbf{w} \in \Fq^{m-n}$ (known only to prover)
\end{itemize}
\item A set of polynomial constraints $\{C_i(\mathbf{v}) = 0\}_{i=1}^{\ell}$ where each $C_i$ is a polynomial over $\Fq$
\end{itemize}

A valid computation satisfies all constraints when variables are assigned appropriate values.

\textbf{Witnesses and Hints:} In constraint systems, complex operations (e.g., field division, polynomial quotients) can be computed efficiently using \textit{witnesses} or \textit{hints}---auxiliary values provided by the prover but verified by constraints. For instance, to verify $a/b = c$:
\begin{enumerate}
\item Prover provides $c$ as a witness
\item Verifier checks the constraint $b \cdot c = a$
\end{enumerate}

This approach transforms expensive operations into cheap verifications, fundamentally altering the cost model.

\textbf{Examples of Constraint Systems:}
\begin{itemize}
\item \textbf{R1CS} (Rank-1 Constraint Systems)\cite{r1cs13}: Constraints of the form $(A\mathbf{v}) \circ (B\mathbf{v}) = (C\mathbf{v})$ where $\circ$ is element-wise multiplication
\item \textbf{Plonkish}\cite{plnk19}: Gate-based constraints with custom gates and copy constraints
\item \textbf{Cairo/AIR}\cite{cairo}: Algebraic Intermediate Representation with polynomial constraints over execution traces
\end{itemize}

\subsubsection{Cost Model}

We measure computational cost in terms of \textit{constraints} (denoted by \textbf{C}), representing the fundamental unit of work in constraint systems. One constraint typically corresponds to one multiplication gate or one algebraic relation to verify.

\textbf{Basic Field Operations in $\Fq$:}
\begin{itemize}
\item \textbf{Addition/Subtraction:} 0C (constraints are linear, free to compose)
\item \textbf{Multiplication:} 1C (one constraint)
\item \textbf{Constant multiplication:} 0C (absorbed into constraint coefficients)
\item \textbf{Division/Inversion:} 2C (one multiplication constraint + one equality check)
\end{itemize}

\textbf{Extension Field Operations:} For $\Fe{k}$ elements represented as $k$ coefficients in $\Fq$:
\begin{itemize}
\item \textbf{Addition:} 0C (component-wise, linear)
\item \textbf{Multiplication:} Traditional approach requires $O(k^2)$ base field multiplications, but our polynomial verification reduces this significantly
\item \textbf{Evaluation at point $x$:} $k \cdot 1C = kC$ (Horner's method: $k-1$ multiplications + $k-1$ additions)
\end{itemize}

\textbf{Commitment Overhead:} Beyond constraint costs, we account for \textit{calldata} or \textit{commitment overhead}---the number of field elements the prover must commit to (via Fiat-Shamir hashing in non-interactive proofs). For blockchain deployments, this overhead directly impacts transaction costs and is often the dominant factor.

\subsubsection{Cost Model Comparison}

Table \ref{tab:cost_comparison} contrasts costs between native implementations and constraint systems:

\begin{table}[h]
\centering
\caption{Cost comparison: Native vs Constraint-based computation}
\label{tab:cost_comparison}
\begin{tabular}{|l|c|c|c|}
\hline
\textbf{Operation} & \textbf{Native Cost} & \textbf{Constraint Cost} & \textbf{Ratio} \\
\hline
$\Fq$ Multiplication & 1$\times$ & 1C & 1$\times$ \\
$\Fq$ Inversion & 100$\times$ & 2C & $\approx$50$\times$ cheaper \\
$\Fe{12}$ Multiplication & 50$\times$ & 54C (traditional) & Similar \\
$\Fe{12}$ Multiplication & 50$\times$ & 39C (our method) & Better \\
\hline
\end{tabular}
\end{table}

\textbf{Key Insight:} Operations that are traditionally expensive (inversions, divisions) become cheap in constraint systems when computed via witnesses. This inverts traditional optimization priorities: we can afford inversions to save multiplications, enabling affine coordinate systems and other techniques impractical in native implementations.

Recognizing and exploiting these unique computational properties is essential for unlocking hidden performance potential. Youssef El Housni's pioneering work \textbf{Pairings in Rank-1 Constraint Systems}~\cite{PR1CS22} represents, to our knowledge, the first systematic exploration of these optimization opportunities, providing a breakthrough in the practicality of recursive proving by thoughtfully leveraging the distinctive cost model inherent to constraint-based systems.
\subsection{Related Work}

The landscape of efficient pairing verification in constraint systems has evolved through several key contributions. El Housni's \textbf{Pairings in Rank-1 Constraint Systems}~\cite{PR1CS22} established the theoretical framework for optimizing pairings in Rank-1 Constraint Systems, introducing techniques for efficient Miller loop implementation and final exponentiation that exploit the unique properties of constraint-based computation models. This work demonstrated that operations traditionally considered expensive in native implementations could be restructured to achieve remarkable efficiency gains in arithmetized environments.

Feltroidprime's subsequent research on \textbf{enhanced extension field multiplication techniques}~\cite{FXFM23} introduced polynomial verification methods using the Schwartz-Zippel lemma for $\Fe{12}$ arithmetic operations. By representing extension field elements as polynomials and verifying their operations through polynomial arithmetic, this approach achieved significant performance improvements in the pairing verification.

Novakovic and Eagen's work \textbf{On Proving Pairings}~\cite{OPP24} presented breakthrough optimization strategies that enable the elimination of final exponentiation overhead through residue witness techniques. Their approach demonstrates that two elements in $\Fe{12}$ can be proven equivalent without expensive final exponentiation by introducing witness elements that satisfy specific algebraic relationships, and effectively embedding the final exponentiation verification into the Miller loop computation itself.

\newcommand{\InlineFor}[2]{\State \textbf{for} #1 \textbf{do} #2}

\section{Pairing Operations as Polynomials}

Algorithm~\ref{alg:pairing} gives the pairing computation combining El Housni's R1CS optimizations~\cite{PR1CS22} with Novakovic-Eagen's residue witness technique~\cite{OPP24}. The 27th root of unity $W$ scales the Miller loop output $f$ to be a cubic residue (see~\cite{OPP24} §4.3); if $f$ is not already a cubic residue, multiplying by $W^w$, $w\in\{1,2\}$ corrects this.

Our contribution is to express each step of this algorithm as a single multi-operand polynomial product verified via the Schwartz-Zippel lemma, rather than as a sequence of individual binary multiplications as in~\cite{FXFM23}.

\begin{algorithm}
    \caption{Optimal Ate Multi-Pairing with Residue witness check}
    \label{alg:multipairing}
    \begin{algorithmic}[1]
    \Require Vectors $\vec{P} = (P_1, \dots, P_N) \in \G_1^N$, $\vec{Q} = (Q_1, \dots, Q_N) \in \G_2^N$, $c \in \Fe{12}, w \in \{0,1,2\}$
    \Statex $\text{Internal } W, W^2 \in \Fe{12}$ \Comment{$W$ is 27th root of unity}
    \Ensure $\prod_{k=1}^N a_{opt}(Q_k, P_k)$
    \State Write $s = 6t + 2 = \sum_{i=0}^{L-1} s_i 2^i$ where $s_i \in \{-1, 0, 1\}$ and $s_L = 1$
    \State $c^{-1} \leftarrow 1 / c$ \Comment{Inverse Residue witness}
    \State $f \leftarrow c^{-1}$ \Comment{Initialize with inverse residue witness}

    \InlineFor {$k = 1$ to $N$} {$T_k \leftarrow Q_k$}

    \Statex

    \For{$i = L - 2$ to $0$}
        \State $f \leftarrow f^2$ \label{pr:ml:sq}
        \State $f \leftarrow f \cdot \prod_{k=1}^N l_{T_k,T_k}(P_k)$ \label{pr:ml:ldbl}

        \InlineFor {$k = 1$ to $N$} {$T_k \leftarrow [2]T_k$}

        \If{$s_i = 1$}
            \State $f \leftarrow f \cdot c^{-1}$ \Comment{Residue witness}\label{pr:ml:rw1bit}
            \State $f \leftarrow f \cdot \prod_{k=1}^N l_{T_k,Q_k}(P_k)$ \label{pr:ml:ladd1bit}
            \InlineFor {$k = 1$ to $N$} {$T_k \leftarrow T_k + Q_k$}

        \ElsIf{$s_i = -1$}
            \State $f \leftarrow f \cdot c$ \Comment{Residue witness}\label{pr:ml:rw-1bit}
            \State $f \leftarrow f \cdot \prod_{k=1}^N l_{T_k,-Q_k}(P_k)$ \label{pr:ml:ladd-1bit}
            \InlineFor {$k = 1$ to $N$} {$T_k \leftarrow T_k - Q_k$}

        \EndIf
    \EndFor
    \Statex

    \InlineFor {$k = 1$ to $N$} {$Q_{1,k} \leftarrow \pi_p(Q_k)$, $Q_{2,k} \leftarrow \pi_p^2(Q_k)$}

    \State $f \leftarrow f \cdot \prod_{k=1}^N l_{T_k,Q_{1,k}}(P_k)$ \label{pr:ml:pi1}

    \InlineFor {$k = 1$ to $N$} {$T_k \leftarrow T_k + Q_{1,k}$}

    \State $f \leftarrow f \cdot \prod_{k=1}^N l_{T_k,-Q_{2,k}}(P_k)$ \label{pr:ml:pi2}

    \Statex
    \If{$w = 1$}
        \State $f \leftarrow f \cdot W$ \label{pr:ml:27r}
    \ElsIf{$w = 2$}
        \State $f \leftarrow f \cdot W^2$ \label{pr:ml:27r2}
    \EndIf

    \State $f \leftarrow f \cdot c^{-q} \cdot c^{q^2} \cdot c^{-q^3}$ \Comment{Uses Frobenius maps, Negative exponents use $c^{-1}$}\label{pr:ml:fe}

    \Statex

    \State \Return $f$
    \end{algorithmic}
\end{algorithm}

\subsection{Comparative Analysis~\label{sec:ops:comp-analysis}}

We compare the two schemes for a 3-pair Miller loop zero-bit step (the most common case in BN254's NAF\cite{naf} representation of $s = 6t+2$). Here $F \in \Fe{12}$ is the accumulator, $L_i \in \text{Sparse}_{01379}$ are the three doubling line functions, $R \in \Fe{12}$ is the updated accumulator, and $Q$ is the quotient polynomial. $\text{Sparse}_{01379}$ denotes the subspace of $\Fe{12}$ elements whose only nonzero coefficients are at degrees $0,1,3,7,9$---the form that BN254 line functions naturally take~\cite{PR1CS22}, reducing their evaluation cost from 12C to 4C.

Adapting~\cite{FXFM23} to lines \ref{pr:ml:sq}--\ref{pr:ml:ldbl} of Algorithm~\ref{alg:pairing} chains four binary multiplications, each producing an intermediate remainder that must be committed:

\begin{align*}
    F(x) \cdot F(x) &\stackrel{?}{=} Q_1(x) \cdot P_{12}(x) + T_1(x) \\
    T_1(x) \cdot L_1(x) &\stackrel{?}{=} Q_2(x) \cdot P_{12}(x) + T_2(x) \\
    T_2(x) \cdot L_2(x) &\stackrel{?}{=} Q_3(x) \cdot P_{12}(x) + T_3(x) \\
    T_3(x) \cdot L_3(x) &\stackrel{?}{=} Q(x) \cdot P_{12}(x) + R(x)
\end{align*}

Each of the five polynomial evaluations ($F$, $T_1$--$T_3$, $R$) costs 12C and each line costs 4C, with all four remainders requiring commitment: \textbf{76C and 48 $\Fq$ commitments}.

Our scheme collapses this into a single identity:\label{sec:op:comp:our}

$$F(x)^2 \cdot L_1(x) \cdot L_2(x) \cdot L_3(x) \stackrel{?}{=} Q(x) \cdot P_{12}(x) + R(x)$$

Only $R$ requires commitment, and only $F$ and $R$ need full 12C evaluation; the three sparse lines cost 4C each: \textbf{41C and 12 $\Fq$ commitments} --- a \textbf{46\% constraint reduction} and \textbf{75\% fewer commitments}. The comparison across all three schemes is summarised in Table~\ref{tab:miller_loop_comparison}.

\begin{table}[h]
\centering
\caption{3-pair Miller loop zero-bit step comparison}
\label{tab:miller_loop_comparison}
\begin{tabular}{|l|c|c|c|}
\hline
\textbf{Scheme} & \textbf{Constraints} & \textbf{$\Fq$ Commitments} & \textbf{vs baseline} \\
\hline
\cite{PR1CS22} & 132C & --- & --- \\
\hline
\cite{FXFM23} & 76C & 48 & $-$42\% \\
\hline
\textbf{Ours} & \textbf{41C} & \textbf{12} & $-$\textbf{69\%} \\
\hline
\multicolumn{4}{|l|}{\footnotesize \cite{PR1CS22}: $36\text{C} + 3\times30\text{C} = 132\text{C}$} \\
\multicolumn{4}{|l|}{\footnotesize \cite{FXFM23}: $5\times12\text{C} + 3\times4\text{C} + 4\text{C} = 76\text{C}$} \\
\multicolumn{4}{|l|}{\footnotesize Ours: $2\times12\text{C} + 3\times4\text{C} + 5\text{C} = 41\text{C}$\quad (excl.\ Fiat-Shamir overhead)} \\
\hline
\end{tabular}
\end{table}


\subsection{Operation Polynomials}\label{sec:op-polys}

We now state the verification equation for each step of Algorithm~\ref{alg:pairing} in a 3-pairing configuration (as used in Groth16 verification). All equations share the form $\text{LHS} = R(x) + Q(x) \cdot P_{12}(x)$; the quotient $Q$ is amortized across all steps via batching (Section~\ref{sec:iop:batching}).

\subsubsection{Zero-bit step}\label{sec:op:ml-0bit}

Lines~\ref{pr:ml:sq}--\ref{pr:ml:ldbl}: one squaring and three doubling lines.

\begin{equation}\label{eq:op:0b}
F^2(x) \cdot L_1(x) \cdot L_2(x) \cdot L_3(x) = R(x) + Q(x) \cdot P_{12}(x)
\end{equation}

$F, R \in \Fe{12}$ (12 coefficients each); $L_1, L_2, L_3 \in \text{Sparse}_{01379}$ (4 coefficients each). \textbf{Cost: 41C, 12 commitments.}

\subsubsection{Non-zero bit step}\label{sec:op:ml-nzbit}

Lines~\ref{pr:ml:rw1bit}--\ref{pr:ml:ladd-1bit}: squaring, residue witness, and both doubling and addition lines for each pair.

\begin{equation}\label{eq:op:n0b}
F^2(x) \cdot W^{\pm1}(x) \cdot L_{\text{d1}}(x) \cdot L_{\text{d2}}(x) \cdot L_{\text{d3}}(x) \cdot L_{\text{a1}}(x) \cdot L_{\text{a2}}(x) \cdot L_{\text{a3}}(x) = R(x) + Q(x) \cdot P_{12}(x)
\end{equation}

$W^{-1}$ is used for negative bits; the equation structure is identical. $W \in \Fe{12}$ encapsulates the equivalence class relationship from~\cite{OPP24}. \textbf{Cost: 69C, 12 commitments.}

\subsubsection{Frobenius correction step}\label{sec:op:ml-frob}

Lines~\ref{pr:ml:pi1}--\ref{pr:ml:pi2}: no squaring; two sparse correction lines per pair for the $\pi_p$ and $-\pi_p^2$ mappings.

\begin{equation}\label{eq:op:fm}
F(x) \cdot L_{1\pi}(x) \cdot L_{2\pi}(x) \cdot L_{3\pi}(x) \cdot L_{1\pi^2}(x) \cdot L_{2\pi^2}(x) \cdot L_{3\pi^2}(x) = R(x) + Q(x) \cdot P_{12}(x)
\end{equation}

\textbf{Cost: 55C, 12 commitments.}

\subsubsection{Final exponentiation step}\label{sec:op:finx}

Lines~\ref{pr:ml:27r}--\ref{pr:ml:fe}: cubic scaling and the Frobenius component ($q - q^2 + q^3$) of the final exponentiation, following~\cite{OPP24}. $R = 1$ for a valid pairing.

\begin{equation}\label{eq:op:fx}
F(x) \cdot c^{-q}(x) \cdot c^{q^2}(x) \cdot c^{-q^3}(x) \cdot W^w(x) = 1 + Q(x) \cdot P_{12}(x)
\end{equation}

$w \in \{0,1,2\}$ selects the cubic root of unity scaling; $W^w$ is sparse with 2 nonzero coefficients. The Frobenius images of $c$ are provided as prover inputs. \textbf{Cost: 56C, 12 commitments.}

\subsection{Security}\label{sec:op-sec}

Each verification equation takes the form $P = \prod_i A_i(x) - Q(x) \cdot P_{12}(x) - R(x)$ where $\deg(P) < 2^8$.\footnote{For $k$-pair pairings, $\deg(P) \approx 88 + 18(k-3)$, well below $2^8$ for all practical $k$. This bound can be extended to $2^{10}$ with negligible soundness impact.} By the Schwartz-Zippel lemma (Section~\ref{sec:bg:sz}):

\textbf{Soundness:} $\delta_s \leq d/|\Fq| < 2^8/2^{254} = 2^{-246}$ for BN254.

\textbf{Completeness:} An honest prover always produces $P(x) = 0$ identically, so the verifier always accepts. Euclidean division guarantees unique $Q, R$ with $\deg(R) < \deg(P_{12}) = 12$.
\section{Full Pairing Proof with Quotient Batching}
\label{sec:iop}

We now combine the polynomial equations from Section~\ref{sec:op-polys} into a complete two-round public-coin interactive proof, further reducing verification cost by batching all quotients into a single accumulated polynomial.

\subsection{Batching Polynomial Identity Testing}\label{sec:iop:batching}

All $n$ polynomial equations from Section~\ref{sec:op-polys} are batched into a single verification. The prover commits to remainders $\{R_i\}_{i=0}^{n-1}$ and receives challenge $z$. Using $z$, the prover forms the accumulated quotient $Q_\text{acc} = \sum_{i=0}^{n-1} z^i \cdot Q_i$. The verifier then samples a second challenge point $c$ and checks:

\begin{equation}
    \label{eq:iop-batched}
    \sum_{i=0}^{n - 1}z^i \cdot [A_i(c) \cdot B_i(c) \cdots X_i(c) - R_i(c)] = Q_\text{acc}(c) \cdot P_{12}(c)
\end{equation}

This reduces $n$ polynomial verifications (degrees 38--76, see Tables~\ref{tbl:op:ml-0bit}--\ref{tbl:op:finx}) to a single evaluation of a degree-76 polynomial.

\subsection{Prover's Algorithms}

The prover computes quotients $Q_i$ and remainders $R_i$ for each equation in Section~\ref{sec:op-polys}. For each step, it calls \textsc{PolyMulDiv}, which multiplies the input polynomials and performs standard Euclidean division by the irreducible $P_{12}$. After receiving challenge $z$, the prover forms $Q_\text{acc} = \sum_{i} z^i Q_i$ and sends $\mathcal{R} = \{R_i\}$ and $Q_\text{acc}$ to the verifier. The full prover procedure is given in Algorithm~\ref{alg:2r:prover}.

% \begin{algorithm}
% \caption{PolyMulDiv: Multiply polynomials and divide by $P_{12}$}
% \label{alg:polymuldiv}
% \begin{algorithmic}[1]
% \Require $P_1, \ldots, P_n \in \Fq[x]$; Divisor $P_{12} \in \Fq[x]$
% \Ensure $Q, R$ s.t.\ $\prod P_i = Q \cdot P_{12} + R$, $\deg(R) < \deg(P_{12})$
% \State $R \leftarrow \prod_{i=1}^{n} P_i$ \gComm{Polynomial product}
% \State Perform Euclidean division of $R$ by $P_{12}$ to get $(Q, R)$
% \State \Return $(Q, R)$
% \end{algorithmic}
% \end{algorithm}

\begin{algorithm}
\caption{Pairing Prover: Preparing polynomials}
\label{alg:2r:prover}
\begin{algorithmic}[1]
\Require $P \in \G_1, Q \in \G_2, c \in \Fe{12}, w \in \{0,1,2\}$
\Statex Internal: $W, W^2 \in \Fe{12}$ \Comment{$W$ is 27th root of unity}
\Ensure Arrays $\mathcal{Q}, \mathcal{R}$ of quotients and remainders
\Statex Write $s = 6t + 2 = \sum_{i=0}^{L-1} s_i 2^i$ where $s_i \in \{-1, 0, 1\}$ and $s_L = 1$
\State $c^{-1} \leftarrow 1 / c$ \Comment{Inverse Residue witness}
\State $f \leftarrow c^{-1}$ \Comment{Initialize with inverse residue witness}
\State $T \leftarrow Q$ \gComm{$T_i$ for each pair for multi-pairing}
\color{gray}
\Statex \small{Every operation involving $T, P \text{ or } Q$ is repeated for $T_i, P_i \text{ and } Q_i$ for multi-pairing.}
\color{black}

\Statex

\State \textbf{function} \texttt{PolyMulDiv} ({$R, P_1, \ldots, P_n \in \Fq[x]$}): \Return $Q, R \in \Fq[x]$ 
\State \qquad s.t.\ $\prod P_i = Q \cdot P_{12} + R$, $\deg(R) < \deg(P_{12})$

\Statex
\Statex \gText{Miller loop}
\For{$i = L - 2$ to $0$}
    \State $L_\text{dbl} \leftarrow l_{T,T}(P)$, $T \leftarrow [2]T$
    \If{$s_i = 0$}
		\State $Q, R \leftarrow {\textit{PolyMulDiv}}(f, f, L_\text{dbl})$ \Comment{add $L_{*i}$ for each $P_i, Q_i$}
    \ElsIf{$s_i = 1$}
        \State $L_\text{add} \leftarrow l_{T,Q}(P)$, $T \leftarrow T + Q$
        \State $Q, R \leftarrow {\textit{PolyMulDiv}}(f, f, c^{-1}, L_\text{dbl}, L_\text{add})$ \Comment{add $L_{*i}$ for each $P_i, Q_i$}
    \ElsIf{$s_i = -1$}
        \State $L_\text{add} \leftarrow l_{T,-Q}(P)$, $T \leftarrow T - Q$
        \State $Q, R \leftarrow {\textit{PolyMulDiv}}(f, f, c, L_\text{dbl}, L_\text{add})$ \Comment{add $L_{*i}$ for each $P_i, Q_i$}
    \EndIf
    \State $f \leftarrow R$, $\mathcal{Q}$.\textit{append}($Q$), $\mathcal{R}$.\textit{append}($R$)
\EndFor
\Statex


\State $Q_1 \leftarrow \pi_p(Q), Q_2 \leftarrow \pi_p^2(Q)$
\State $L_\pi \leftarrow l_{T,Q_1}(P)$
\State $T \leftarrow T + Q_1$
\State $L_{-\pi^2} \leftarrow l_{T,-Q_2}(P)$
\State $Q, R \leftarrow {\textit{PolyMulDiv}}(f, L_\pi, L_{-\pi^2})$ \Comment{add $L_{*i}$ for each $P_i, Q_i$}
\State $f \leftarrow R$, $\mathcal{Q}$.\textit{append}($Q$), $\mathcal{R}$.\textit{append}($R$)

\Statex

\Statex \gText{Final exponentiation: Cubic scaling + Frobenius mappings}

\If{$w = 0$}
    \State $Q, R \leftarrow {\textit{PolyMulDiv}}(f, c^{-q}, c^{q^2}, c^{-q^3})$
\ElsIf{$w = 1$}
    \State $Q, R \leftarrow {\textit{PolyMulDiv}}(f, W, c^{-q}, c^{q^2}, c^{-q^3})$
\ElsIf{$w = 2$}
    \State $Q, R \leftarrow {\textit{PolyMulDiv}}(f, W^2, c^{-q}, c^{q^2}, c^{-q^3})$
\EndIf
\State $\mathcal{Q}$.\textit{append}($Q$), $\mathcal{R}$.\textit{append}($R$) \gComm{Final $\mathcal{Q}$ and $\mathcal{R}$ accumulation, skip $f$}

\Statex

\State \Return $\mathcal{Q}, \mathcal{R}$
\end{algorithmic}
\end{algorithm}

\subsection{Verifier's Algorithms}

The verifier receives $P \in \G_1$, $Q \in \G_2$, residue witness $c$~\cite{OPP24}, the $W$ (27th root of unity), the remainder array $\mathcal{R}$, and $Q_\text{acc}$ from the prover. Using challenge $z$ from the first round and a freshly sampled $x$, it mirrors the prover's pairing loop while evaluating each polynomial at $x$, accumulating the RLC $e_\text{acc}$. The final check is Equation~\ref{eq:iop-batched}. The full verifier is given in Algorithm~\ref{alg:2r:verifier}.

% \begin{algorithm}
% \caption{EvalPoly: Evaluate polynomial product minus remainder via Horner's method}
% \label{alg:evalpolyprod}
% \begin{algorithmic}[1]
% \Require Polynomials $P_1, \ldots, P_n,\, R \in \Fq[x]$ of degree $\leq d$, evaluation point $z \in \Fq$
% \Ensure $e \in \Fq$ where $e = \prod_{i=1}^{n} P_i(z) - R(z)$

% \State \textbf{function} \texttt{EvalPoly} ({$R, P_1, \ldots, P_n \in \Fq[x]$}): \Return ($\prod_{i=1}^{n} P_i(z) - R(z) \in \Fq$)

% \State \textbf{function} \texttt{Eval} ({$P = (p_0, p_1, \ldots, p_d),\ z$}): \Return ($evaluation \in \Fq$)

% \State \gText{~~~~Uses Horner's method, body not included for brevity}

% \State $e \leftarrow \Call{Eval}{P_1, z}$ \gComm{Evaluate first polynomial}
% \For{$i = 2$ to $n$}
%     \State $e \leftarrow e \cdot \Call{Eval}{P_i, z}$ \gComm{Multiply in each subsequent evaluation}
% \EndFor
% \State $e \leftarrow e - \Call{Eval}{R, z}$ \gComm{Subtract remainder evaluation}
% \State \Return $e$
% \end{algorithmic}
% \end{algorithm}

\begin{algorithm}
\caption{Pairing Verifier: Evaluating remainders}
\label{alg:2r:verifier}
\begin{algorithmic}[1]
\Require $P \in \G_1, Q \in \G_2, c \in \Fe{12}, w \in \{0,1,2\}$, $\mathcal{R} \text{ array of remainders}$, $Q_\text{acc} \in \Fq[x]$
\Statex Internal: 27th root of unity, $W, W^2 \in \Fe{12}$, first round challenge, $z \xleftarrow{\$} \Fq$
\Ensure $pairing(P, Q) \stackrel{?}{=} 1$

\Statex Write $s = 6t + 2 = \sum_{i=0}^{L-1} s_i 2^i$ where $s_i \in \{-1, 0, 1\}$ and $s_L = 1$

\State $c^{-1} \leftarrow 1 / c$ \Comment{Inverse Residue witness}
\State $f \leftarrow c^{-1}$ \Comment{Initialize with inverse residue witness}
\State $T \leftarrow Q$ \gComm{$T_i$ for each pair for multi-pairing}
\State $x \xleftarrow{\$} \Fq$ \gComm{Sample evaluation point from field}
\State $e_\text{acc} \leftarrow 0$, $a_\text{rlc} \leftarrow 1$ \gComm{Evaluation accumulator and random linear combination multiplier}
\color{gray}
\Statex \small{Every operation involving $T, P \text{ or } Q$ is repeated for $T_i, P_i \text{ and } Q_i$ for multi-pairing.}
\color{black}

\Statex

\State \textbf{function} \texttt{EvalPoly} ({$R, P_1, \ldots, P_n \in \Fq[x]$}): \Return ($\prod_{i=1}^{n} P_i(z) - R(z) \in \Fq$)

\Statex
\Statex \gText{Miller loop}
\For{$i = L - 2$ to $0$}
    \State $L_\text{dbl} \leftarrow l_{T,T}(P)$, $T \leftarrow [2]T$
    \State $R \leftarrow \mathcal{R}.\textit{pop\_front}()$ \gComm{Fetch R from prover}
    \If{$s_i = 0$}
		\State $e_\text{acc} \leftarrow e_\text{acc} + a_\text{rlc} * \texttt{EvalPoly}(R, x, f, f, L_{\text{dbl}})$ \Comment{add $L_{*i}$ for each $P_i, Q_i$}
    \ElsIf{$s_i = 1$}
        \State $L_\text{add} \leftarrow l_{T,Q}(P)$, $T \leftarrow T + Q$
        \State $e_\text{acc} \leftarrow e_\text{acc} + a_\text{rlc} * \texttt{EvalPoly}(R, x, f, f, c^{-1}, L_\text{dbl}, L_\text{add})$ \Comment{add $L_{*i}$ for each $P_i, Q_i$}
    \ElsIf{$s_i = -1$}
        \State $L_\text{add} \leftarrow l_{T,-Q}(P)$, $T \leftarrow T - Q$
        \State $e_\text{acc} \leftarrow e_\text{acc} + a_\text{rlc} * \texttt{EvalPoly}(R, x, f, f, c, L_\text{dbl}, L_\text{add})$ \Comment{add $L_{*i}$ for each $P_i, Q_i$}
    \EndIf
    \State $f \leftarrow R$, $a_\text{rlc} \leftarrow a_\text{rlc} \cdot z$ \gComm{Update random linear combination factor}
\EndFor
\Statex

\State $R \leftarrow \mathcal{R}.\textit{pop\_front}()$ \gComm{Fetch R from prover}
\State $Q_1 \leftarrow \pi_p(Q), Q_2 \leftarrow \pi_p^2(Q)$
\State $L_\pi \leftarrow l_{T,Q_1}(P)$
\State $T \leftarrow T + Q_1$
\State $L_{-\pi^2} \leftarrow l_{T,-Q_2}(P)$
\State $e_\text{acc} \leftarrow e_\text{acc} + a_\text{rlc} * \texttt{EvalPoly}(R, x, f, L_\pi, L_{-\pi^2})$ \Comment{add $L_{*i}$ for each $P_i, Q_i$}
\State $f \leftarrow R$, $a_\text{rlc} \leftarrow a_\text{rlc} \cdot z$ \gComm{Update random linear combination factor}

\Statex

\Statex \gText{Final exponentiation: Cubic scaling + Frobenius mappings}
\State $R \leftarrow \mathcal{R}.\textit{pop\_front}()$ \gComm{Fetch R from prover}
\If{$w = 0$}
    \State $e_\text{acc} \leftarrow e_\text{acc} + a_\text{rlc} * \texttt{EvalPoly}(R, x, f, c^{-q}, c^{q^2}, c^{-q^3})$
\ElsIf{$w = 1$}
    \State $e_\text{acc} \leftarrow e_\text{acc} + a_\text{rlc} * \texttt{EvalPoly}(R, x, f, W, c^{-q}, c^{q^2}, c^{-q^3})$
\ElsIf{$w = 2$}
    \State $e_\text{acc} \leftarrow e_\text{acc} + a_\text{rlc} * \texttt{EvalPoly}(R, x, f, W^2, c^{-q}, c^{q^2}, c^{-q^3})$
\EndIf

\Statex

\State \Return $Q_\text{acc}(x) \stackrel{?}{=} e_\text{acc} \land R \stackrel{?}{=} 1$ \gComm{Polynomial identity and expected pairing output check}
\end{algorithmic}
\end{algorithm}

\subsection{Interactive Proof} \label{sec:iop:interactive}

\IOP{Two-Round Pairing IOP}{Verifier}{Prover}{
    \RDoes{Computes pairings with Algorithm~\ref{alg:2r:prover} and prepares\\ remainders and quotients arrays $\mathcal{R}$ and $\mathcal{Q}$}{1}
    \RtoL{Round I: Commits \textit{pairing inputs} and remainders array $\mathcal{R}$}{2.5}
    \LtoR{Challenge $z \xleftarrow{\$} \Fq$}{3.5}
    \RDoes{Computes $Q_{\text{acc}} = \sum_{i=0}^{n-1} z^i \cdot Q_i$ using $\mathcal{Q}$ and $z$}{4.5}
    \RtoL{Round II: Commit $Q_\text{acc}$}{5.5}
    \LDoes{Samples $x \xleftarrow{\$} \Fq$ and runs Algorithm~\ref{alg:2r:verifier} to test\\polynomial identities at $x$, Eq~\ref{eq:iop-batched}}{6.5}
    \LtoR{Accept/Reject}{8}
}

\subsection{Fiat-Shamir Transformation} \label{sec:iop:fs}

Applying the Fiat-Shamir transformation~\cite{fiatshamir} with random oracle $H$ to the two-round protocol of Section~\ref{sec:iop:interactive} gives a non-interactive proof with knowledge soundness $\delta_s + 3(m^2+1)2^{-\lambda}$~\cite{IOP-NI}, where $m=2$ is the number of oracle queries (one per round) and $\lambda$ is the challenge bit-length; Section~\ref{sec:iop:soundness} instantiates this bound.

\subsubsection{Non-Interactive Proof} \label{sec:iop:noninteractive}

The prover runs Algorithm~\ref{alg:2r:prover} to obtain $(\mathcal{Q}, \mathcal{R})$, derives the challenge
$$z = H(\text{\textit{pairing inputs}},\, c,\, W^w,\, \mathcal{R})$$
from the transcript instead of receiving it, computes $Q_{\text{acc}} = \sum_{i=0}^{n-1} z^i \cdot Q_i$, and outputs the proof $(\text{\textit{pairing inputs}},\, c,\, W^w,\, \mathcal{R},\, Q_\text{acc})$. The verifier recomputes $z$, derives the evaluation point $x = H(z,\, Q_\text{acc})$, and runs Algorithm~\ref{alg:2r:verifier}, testing the polynomial identities via Equation~\eqref{eq:iop-batched}.

The soundness bound of Theorem~\ref{thm:soundness} carries over because every prover-chosen value entering Equation~\eqref{eq:iop-batched} is absorbed before the challenge it must not depend on: $\mathcal{R}$ and the hinted witnesses $c, W^w$ before $z$, and $Q_\text{acc}$ with $z$ chained into the second query before $x$.

\subsubsection{Instantiating the Random Oracle}\label{sec:iop:fs:ro}

The protocol is agnostic to how $H$ is realised; what varies across execution environments is the cost of a \emph{commitment}, i.e.\ of absorbing one $\Fq$ element into the transcript. We therefore qualify the commitment counts of Section~\ref{sec:ops:comp-analysis} and Table~\ref{tab:garaga_performance} for our two instantiations.

\paragraph{CairoVM (Poseidon).} \label{sec:iop:fs:cairo}
Prover-supplied hints are passed as transaction calldata and absorbed into a running Poseidon~\cite{poseidon} sponge; $\Fq$ elements are represented by 96-bit limbs, with two limbs packed into each \texttt{Felt252} absorption. Both hashing cost and calldata therefore scale linearly with the number of committed elements, which is why the COMMITS column of Table~\ref{tab:garaga_performance} is reported alongside modular operations: halving the committed elements roughly halves the transcript-related share of VM steps and the hint calldata.

\paragraph{gnark (\texttt{api.Commit}).} \label{sec:iop:fs:gnark}
The two rounds map to the \textbf{two} sequential \texttt{api.Commit} calls in \texttt{performDeferredRingChecks} (Section~\ref{sec:polyring:mapping}): the first commits the limbs of the queued remainders and hinted witnesses, yielding $z$; the second commits $z$ together with the limbs of $Q_\text{acc}$, yielding $x$. The cost of \texttt{Commit} is backend-specific. In Groth16, gnark implements the commitment scheme of Belling et al.~\cite{BellingCCS23,gnarkcommit}: the committed wires are aggregated into a Pedersen-style $\G_1$ commitment, and the challenge is a hash-to-field of the commitment point, computed \emph{natively} by prover and verifier---no in-circuit hashing. Each \texttt{Commit} costs the prover two $\G_1$ MSMs of size equal to the number of committed wires (the commitment and its proof of knowledge); the proof grows by one $\G_1$ point per commitment plus a single proof of knowledge folded across commitments (64 bytes each, uncompressed); and verification gains one two-pairing product check, one hash-to-field, and one extra scalar multiplication and point addition in the public-input MSM.
In PLONK, committed wires reuse the existing wire-commitment machinery, so the marginal overhead per commitment is lower still, but each committed wire adds a constraint.

\paragraph{Impact.} The number of elements committed to determines hashing and calldata in CairoVM and the commitment-MSM size in gnark. Reducing the absolute number of committed elements, by 53\% through our consolidation (Section~\ref{sec:ops:comp-analysis}), directly reduces end-to-end cost regardless of the instantiation, without adding commitment rounds.


\subsection{Security Analysis}

Throughout this section, $z, x \in \mathbb{F}_q$ denote field element challenges (scalars),
while polynomials are elements of $\mathbb{F}_q[X]$ with $X$ as the formal indeterminate.
In particular, $z^j$ is a scalar coefficient, not a polynomial term.

\subsubsection{Algebraic Intuition}\label{sec:iop:alg-intuition}

Before the formal soundness proof it is instructive to consider why an incorrect
remainder is difficult to hide.  Suppose a malicious prover corrupts exactly one remainder:
$R_j' = R_j + \Delta_j$ with $\Delta_j \in \mathbb{F}_q[X]$, $\Delta_j \neq 0$,
$\deg(\Delta_j) \leq 11$.  Since $z$ is obtained \emph{after} the commitment to
$\mathcal{R}$, the prover cannot adjust any other $R_k$ to compensate for $\Delta_j$.
To successfully deceive the verifier, the prover must find a valid forged quotient
$Q'_\text{acc} \in \mathbb{F}_q[X]$ satisfying

$$\sum_{i=0}^{n-1} z^i(A_i - R_i') = Q'_\text{acc} \cdot P_{12}$$

All $R_i' = R_i$ except for $R_j' = R_j + \Delta_j$. Substituting $R_j' = R_j + \Delta_j$ gives
\begin{align*}
Q'_\text{acc} \cdot P_{12} &= Q_\text{acc} \cdot P_{12} - z^j\Delta_j(x) \\
Q'_\text{acc} &= Q_\text{acc} - \frac{z^j\Delta_j(x)}{P_{12}(x)}.
\end{align*}
Here $z^j \in \mathbb{F}_q$ is a nonzero scalar (with overwhelming probability over the
choice of $z$) and $\Delta_j(x) \in \mathbb{F}_q[X]$ is a nonzero polynomial of degree
at most~$11$.  Since $P_{12}$ is irreducible of degree~$12$ and
$\deg(z^j\Delta_j) \leq 11 < 12 = \deg(P_{12})$, the polynomial $P_{12}$ does not
divide $z^j\Delta_j$ in $\mathbb{F}_q[X]$; hence $z^j\Delta_j/P_{12}$ is not a
polynomial. Consequently no valid $Q'_\text{acc} \in \mathbb{F}_q[X]$ can satisfy the
equation for this single-corruption case.


The same degree argument extends immediately to an adversary who corrupts
\emph{multiple} remainders: for any scalar $z$, the weighted sum
$\sum_i z^i \Delta_i(x)$ is still a polynomial in $x$ of degree at most~$11$
(since each $z^i$ is a scalar and each $\deg(\Delta_i) \leq 11$), so it cannot
be divisible by $P_{12}$ of degree~$12$ unless it is the zero polynomial, which
requires every $\Delta_i = 0$.  The formal soundness proof below re-derives this
uniformly via Schwartz-Zippel.


\subsubsection{Soundness}\label{sec:iop:soundness}

\begin{theorem}[Soundness]\label{thm:soundness}
Let $n$ be the number of polynomial equations batched via a random linear combination into the bivariate polynomial identity:
$$P_{\mathrm{final}}(z, x) = \sum_{i=0}^{n-1} z^i \cdot \Bigl[\prod_{j} P_{i,j}(x) - R_i(x)\Bigr] - Q_{\mathrm{acc}}(x) \cdot P_{12}(x) = 0$$
Let the maximum degree in $x$ of the polynomial products be $d_x < 2^8$. For a prover committing to incorrect remainders, the probability of successful forgery $\delta_s$ over randomly chosen $(z,x) \in \mathbb{F}_q^2$ is bounded by:
$$\delta_s \leq \frac{n - 1 + d_x}{|\mathbb{F}_q|}$$
\end{theorem}

\begin{proof}
If the prover commits to an incorrect remainder $R_i(x)$, the evaluation yields a non-zero error polynomial $P_{\mathrm{err}}(z, x)$. The structure of $P_{\mathrm{final}}(z, x)$ establishes the following degree bounds for $P_{\mathrm{err}}(z, x)$:
\begin{itemize}
    \item \textbf{Degree in $z$:} at most $n - 1$ (due to the $z^i$ weighting).
    \item \textbf{Degree in $x$:} at most $d_x < 2^8$ (due to the polynomial products).
\end{itemize}
By the Schwartz-Zippel lemma, the probability that the non-zero polynomial $P_{\mathrm{err}}$ evaluates to zero is bounded by the total degree divided by the field size:
$$\Pr[P_{\mathrm{err}}(z, x) = 0] \leq \frac{\deg_z(P_{\mathrm{err}}) + \deg_x(P_{\mathrm{err}})}{|\mathbb{F}_q|} \leq \frac{n - 1 + d_x}{|\mathbb{F}_q|}$$
\end{proof}

\textbf{Concrete parameters for BN254}: With $n = 67$ equations and $d_x \leq 256$:
$$\delta_s \leq \frac{66 + 256}{2^{254}} < \frac{2^9}{2^{254}} = 2^{-245}.$$

\textbf{Non-interactive proof via Fiat-Shamir} \label{sec:iop:soundness:fs}: From Section~\ref{sec:iop:fs}, under the random oracle model the
knowledge soundness error becomes~\cite{IOP-NI}:
$$\delta_s' = \delta_s + 3(m^2 + 1) \cdot 2^{-\lambda}$$
where $m = 2$ (one query per round) and $\lambda = 254$ (challenge bit-length), giving
$$\delta_s' \leq 2^{-245} + 15 \cdot 2^{-254} < 2^{-244}.$$

\begin{remark}[Correctness of polynomial evaluation]\label{rem:overflow}
Theorem~\ref{thm:soundness} assumes that the verifier evaluates all polynomials at the
challenge point $x$ without arithmetic error.  The circuit implementation uses the
deferred-reduction inner-product optimisation of Section~\ref{sec:polyring:performance}
(Section~\ref{sec:emulated:flat}), which accumulates limb-wise products over $\mathbb{Z}$
before a single final reduction.  Correctness of this optimisation requires that no
intermediate limb value reaches the native field characteristic $p$.  We verify this
bound concretely for BN254 in Lemma~\ref{lem:overflow} (Section~\ref{sec:emulated:flat}).
\end{remark}

\subsubsection{Completeness}

For an honest prover executing Algorithm~\ref{alg:2r:prover} with valid inputs:

\begin{enumerate}
\item \textbf{Polynomial construction}: At each operation, the prover computes following polynomials:
   $$\prod_{j} P_{i,j} = Q_i \cdot P_{12} + R_i$$

\item \textbf{RLC formation}: Given challenge $z$, the prover forms:
   $$Q_{\text{acc}} = \sum_{i=0}^{n-1} z^i \cdot Q_i$$

\item \textbf{Verifier's check}: The verifier evaluates at point $x$:
   \begin{align*}
   \text{LHS} &= \sum_{i=0}^{n-1} z^i \cdot \left[\prod_{j} P_{i,j}(x) - R_i(x)\right]\\
   &= \sum_{i=0}^{n-1} z^i \cdot Q_i(x) \cdot P_{12}(x) \quad \text{(by construction)}\\
   &= P_{12}(x) \cdot \sum_{i=0}^{n-1} z^i \cdot Q_i(x)\\
   &= P_{12}(x) \cdot Q_{\text{acc}}(x) = \text{RHS}
   \end{align*}
\end{enumerate}

Since the equality holds for any evaluation point $x$, the protocol satisfies perfect completeness.

% https://github.com/shramee/cairo_pairing/blob/main/packages/bn254_u256/src/pairing/schzip/steps.cairo

% ============================================================
% Section: Polynomial Ring Toolkit
% ============================================================

\section{Polynomial Ring Toolkit}
\label{sec:polyring}

The protocol of Section~\ref{sec:iop} is packaged as a reusable toolkit in
the \texttt{emulated} package of gnark~\cite{gnaF}, parameterised by any
emulated field $\Fq$ (via \texttt{FieldParams} type $T$) and any modulus
polynomial $P \in \Fq[x]$. The toolkit operates in a recursive proof
setting where the prover's computations are performed outside the circuit
via \texttt{hints} (Section~\ref{sec:polyring:hints}). And verifier's checks
are implemented as constraints inside the circuit.

% ----------------------------------------------------------
\subsection{Ring Representation}
\label{sec:polyring:rep}

A ring element is essentially a coefficient slice \texttt{[]*Element[T]} in
ascending degree order. A \texttt{nil} entry is treated as zero by
\texttt{InnerProductNoReduce}---the term is skipped entirely---so sparse
polynomials carry no wasted constraint variables.

The optional \texttt{modEvalFn} of type
$\texttt{EvalFnType[T]} = ([\texttt{*Element[T]}]) \to \texttt{*Element[T]}$
allows handling small-integer coefficients efficiently by multiplying by the small positive
integer first and negating afterwards, rather than multiplying by the
full-width field representative $q - c$.

\subsubsection*{Example: BN254 $\mathbb{F}_{p^{12}}$}
\label{sec:polyring:rep:sparse}

$P_{12}(x) = x^{12} - 18x^6 + 82$ has nine zero coefficients (stored as
\texttt{nil}) and two small non-zero ones. The \texttt{modEvalFn} computes
$P_{12}(\alpha) = \texttt{MulConst}(\alpha^0, 82) + \texttt{Neg}(\texttt{MulConst}(\alpha^6, 18)) + \alpha^{12}$,
keeping intermediate values small instead of multiplying by $q - 18$.

% ----------------------------------------------------------
\subsection{Hint Functions}
\label{sec:polyring:hints}

Hints allow defining computations outside of a circuit. A hint defines an
annotated function in the circuit at compile time. Inputs and outputs are
injected at the solving time. We define these two hints to implement the 
prover's steps in the two-round protocol of Section~\ref{sec:iop:interactive}.

\paragraph{\texttt{polyRingMulHint}} implements \textsc{PolyMulDiv}
(Algorithm~\ref{alg:polymuldiv}): it multiplies the input polynomials, divides
by $P$ to obtain $(Q, R)$, and serialises the result as
$\texttt{nbQTerms} \mid Q\text{ limbs} \mid \texttt{nbRemTerms} \mid R\text{ limbs}$.
Inputs are packed as $\texttt{nbBits} \mid \texttt{nbLimbs} \mid \texttt{nbPoly} \mid q \mid \text{inputs} \mid P$,
where each polynomial is prefixed by its term count.

\paragraph{\texttt{quotientsRLCHint}} computes the accumulated quotient
$Q_{\mathrm{acc}}[j] = \sum_{i} z^i \cdot Q_i[j] \bmod q$ coefficient-wise,
corresponding to the prover's step in Algorithm~\ref{alg:2r:prover}.

% ----------------------------------------------------------
\subsection{API}
\label{sec:polyring:api}

\subsubsection*{Register a ring group.}

\begin{verbatim}
group := f.NewPolyRingCheck(mod, modEvalFn)
\end{verbatim}

Creates a \texttt{PolyRingGroupChecks[T]}, appends it to the field's deferred
check list, and returns a handle shared by all multiplications under the same
modulus. Pass \texttt{nil} for \texttt{modEvalFn} to use the dense fallback.
Multiple groups with different moduli can coexist in one circuit.


\subsubsection*{Submit a multiplication.}
\begin{verbatim}
rem, err := f.MulPolyRings(group, A, B, C, ...)
\end{verbatim}
Computes $R = (\prod_j A_j) \bmod P$ via \texttt{polyRingMulHint} outside the
circuit, range-checks $R$'s limbs inside (\texttt{packLimbs} with
\texttt{strict=true}), and registers the check $(A,B,C;\,R;\,Q)$ in
\texttt{group}. The quotient $Q$ is not range-checked; it is consumed only by
the RLC hint outside the circuit. The returned \texttt{rem} is a fully reduced
\texttt{*Poly[T]} and can be fed directly into the next \texttt{MulPolyRings}
call.

\subsubsection*{Deferred verification.}
\begin{verbatim}
f.performDeferredRingChecks(api)    // called once at circuit finalisation
\end{verbatim}
Implements Equation~\eqref{eq:iop-batched} via two sequential \texttt{api.Commit}
calls~\cite{gnarkcommit} (producing challenges $z$ and $x$) and one \texttt{AssertIsEqual} per
group. Gnark derives $z = \mathsf{MiMC}(\mathsf{root})$ from the first call and $x = \mathsf{MiMC}(z)$
from the second, chaining the challenges across rounds (see Section~\ref{sec:iop:noninteractive}).
No in-circuit work beyond remainder range-checks is performed until this point.

% ----------------------------------------------------------
\subsection{Usage in Practice}
\label{sec:polyring:usage}

Register the ring group once at initialisation:

\begin{verbatim}
fp, _ := emulated.NewField[emulated.BN254Fp](api)
modPoly := fp.MakePoly([]any{82, 0, 0, 0, 0, 0, -18, 0, 0, 0, 0, 0, 1})
ring := fp.NewPolyRingCheck(modPoly, nil)
\end{verbatim}

\texttt{MakePoly} accepts Go integer literals and maps zeros to \texttt{nil},
so the slice is sparse from construction.

Each Miller-loop step then calls:
\begin{verbatim}
rem, _ = f.MulPolyRings(ring, f, f, L1, L2, L3)
\end{verbatim}
and \texttt{performDeferredRingChecks} is invoked once by the gnark backend at
finalisation, running the full two-round protocol of Section~\ref{sec:iop}.

% ----------------------------------------------------------
\subsection{Protocol and Code Flow}
\label{sec:polyring:flow}

The following paragraphs trace each step of the two-round protocol
(Section~\ref{sec:iop:interactive}) to the concrete function that implements it.

The verifier is essentially the circuit which encodes all the operations into constraints which can be later verified. The prover is essentially the computation done outside the circuit via \texttt{hints}.

\subsubsection*{Prover computes $Q_i, R_i$ (Algorithm~\ref{alg:2r:prover})}
%
Each \texttt{MulPolyRings} calls \texttt{polyRingMulHint}, results are provided on call, $Q, R$
for the operation stored.

\subsubsection*{Round~I: commit to $\{R_i\}$, receive challenge $z$}
%
\texttt{performDeferredRingChecks} collects all $R_i$ limbs across all groups
and calls \texttt{api.Commit} to obtain $z$ in the circuit. Internally, gnark folds
all committed limbs into a global linear combination $\mathsf{root}$ and derives
$z = \mathsf{MiMC}(\mathsf{root})$~\cite{gnarkcommit}.

\subsubsection*{Prover forms $Q_{\mathrm{acc}} = \sum_i z^i Q_i$ (Algorithm~\ref{alg:2r:prover})}
%
$z$ from the circuit is passed to \texttt{quotientsRLCHint}, accumulating quotients coefficient-wise in $\bmod\,q$ arithmetic.

\subsubsection*{Round~II: commit to $Q_{\mathrm{acc}}$, receive challenge $x$}
%
The $Q_{\mathrm{acc}}$ limbs are passed to a second \texttt{api.Commit} call.
Gnark chains the hash: $x = \mathsf{MiMC}(z)$, binding $x$ to the full transcript
including all $R_i$ and $Q_{\mathrm{acc}}$~\cite{gnarkcommit}.

\subsubsection*{Verifier evaluates each polynomial at $x$ (Algorithm~\ref{alg:evalpolyprod})}
%
\texttt{evalPolyWithChallenge} calls \texttt{InnerProductNoReduce} over the
pre-computed power vector $\{x^i\}$; the result is memoised in
\texttt{Poly.evaluation} to avoid re-evaluation.

\paragraph{Verifier checks Equation~\eqref{eq:iop-batched}: $\sum_i z^i(\prod_j A_{i,j}(x) - R_i(x)) = Q_{\mathrm{acc}}(x)\cdot P(x)$.}
\texttt{InnerProductNoReduce} on the per-check defects weighted by $\{z^i\}$
gives \texttt{lhs}; \texttt{MulNoReduce} of $Q_{\mathrm{acc}}(x)$ and $P(x)$
(via \texttt{modEvalFn} or the dense fallback) gives \texttt{rhs};
\texttt{AssertIsEqual(lhs, rhs)} closes the circuit.


% ----------------------------------------------------------
\subsection{Implementation Details and Performance}
\label{sec:polyring:performance}

\subsubsection{Schoolbook vs Horner's Method for Polynomial Evaluation}
\label{sec:emulated:flat}

Horner's method for evaluations across emulated elements causes a linear increase in the limbs count per coefficient.
This necessitates reductions after each multiplication.
Instead, we cache powers of the evaluation point $x$ in $\mathcal{X} = (x^0, x^1, \ldots, x^d)$ and evaluate each polynomial as an inner product of $\mathcal{X}$ and the coefficient array $\mathbf{F}$:
$$\mathcal{X} \cdot \mathbf{F} = \mathcal{X}_0 \cdot F_0 + \mathcal{X}_1 \cdot F_1 + \cdots + \mathcal{X}_d \cdot F_d$$

In the emulated field context, \textsc{Eval} in Algorithm~\ref{alg:evalpolyprod} is replaced by this inner product with $\mathcal{X}$. This allows us to use the inner product optimization described below.

\subsubsection{Native field operations for inner products}

Modular reductions can be deferred or entirely elided for operations expressible as inner products, such as polynomial evaluations and random linear combinations. This optimization leverages the capacity headroom provided by representing small emulated limbs (e.g., 64-bit limbs) within a significantly larger native field element.

\paragraph{Deferred Reduction:} Intermediate scalars are permitted to grow over the integers $\mathbb{Z}$ during limb-wise multiplications and linear accumulations.

\paragraph{Preservation of Identity:} Algebraic identities are strictly preserved provided the maximum accumulated limb magnitude remains bounded below the native characteristic $p$.

\paragraph{Handling Asymmetry:} In circuits with asymmetric multiplicative depths, operands may accumulate as disparate polynomial multiples of the emulated characteristic $p_{\text{emul}}$. Consequently, a terminal single-element reduction suffices for equivalence testing, bypassing the constraint overhead of intermediate normalization.
\section{Conclusion}

\subsection{Target Environments}

Our approach delivers maximum impact in constraint-based environments where inversions are cheap and calldata is expensive. \textbf{Arithmetized systems} (Cairo VM, R1CS, Plonkish) benefit because Schwartz-Zippel evaluation is efficient and field inversions cost only 2C. \textbf{Calldata-constrained blockchains} (Ethereum L2s, BitVM) benefit most directly: the 77\% commitment reduction means more proofs fit within transaction size limits. \textbf{Recursive proof systems} benefit because the savings compound at each recursion level. \textbf{ZKVMs} (RISC Zero, SP1, Jolt) benefit naturally as they express computation as constraint systems compatible with our polynomial framework.

\subsection{Implementation and Performance}

Table~\ref{tab:main_performance} presents measurements from Garaga's production Cairo implementation on Starknet. Our theoretical analysis (Section~\ref{sec:op:comp:our}) predicted 46\% constraint reduction and 75\% commitment reduction per zero-bit step. The measured end-to-end results for 3-pair BN254 Miller loop are: \textbf{45\% MULMOD reduction} (14,456$\to$7,975), \textbf{77\% POSEIDON reduction} (3,758$\to$876), and \textbf{53\% VM cycle reduction}.

The hash reduction slightly exceeds the per-step prediction because non-zero bits save even more commitments (two line functions instead of one per pair). The commitment savings also scale with $n$: traditional schemes require $O(k \cdot n)$ Fiat-Shamir commitments; our scheme requires only $O(k)$, independent of the number of pairs.

\begin{table*}[t]
\centering
\caption{Performance comparison: Cairo implementation on Starknet}
\label{tab:main_performance}
\begin{tabular}{|l|r|r|r|r|r|r|r|}
\hline
\textbf{Operation} & \multicolumn{2}{c|}{\textbf{MULMOD}} & \multicolumn{2}{c|}{\textbf{ADDMOD}} & \multicolumn{2}{c|}{\textbf{POSEIDON}} & \textbf{Speedup} \\
\cline{2-7}
& Before & After & Before & After & Before & After & \\
\hline
\multicolumn{8}{|c|}{\textbf{Miller Loop (BN254)}} \\
\hline
Single pair & 5,984 & 3,303 & 5,927 & 3,228 & 1,810 & 828 & 47.7\% \\
% n=2 pairs & 10,132 & 5,639 & 10,107 & 5,576 & 2,740 & 852 & 51.1\% \\
Three pairs & 14,456 & 7,975 & 14,463 & 7,924 & 3,758 & 876 & 53.2\% \\
\hline
\textbf{Average} & \multicolumn{2}{c|}{\textbf{44.8\%}} & \multicolumn{2}{c|}{\textbf{45.2\%}} & \multicolumn{2}{c|}{\textbf{76.7\%}} & \textbf{50.6\%} \\
\hline
\hline
\multicolumn{8}{|c|}{\textbf{Complete Pairing (BN254)}} \\
\hline
Single pair & 10,670 & 7,984 & 13,150 & 10,446 & 3,741 & 2,759 & 23.8\% \\
% n=2 pairs & 14,818 & 10,320 & 17,330 & 12,794 & 4,671 & 2,783 & 31.7\% \\
Three pairs & 19,142 & 12,656 & 21,686 & 15,142 & 5,689 & 2,807 & 37.1\% \\
\hline
\textbf{Average} & \multicolumn{2}{c|}{\textbf{33.9\%}} & \multicolumn{2}{c|}{\textbf{30.2\%}} & \multicolumn{2}{c|}{\textbf{50.7\%}} & \textbf{30.9\%} \\
\hline
\hline
\multicolumn{8}{|c|}{\textbf{Miller Loop (BLS12\_381)}} \\
\hline
Single pair & 4,936 & 2,672 & 4,966 & 2,686 & 1,580 & 790 & 47.2\% \\
% n=2 pairs & 8,030 & 4,418 & 8,171 & 4,525 & 2,276 & 812 & 50.9\% \\
Three pairs & 11,356 & 6,164 & 11,608 & 6,364 & 3,088 & 834 & 53.8\% \\
\hline
\textbf{Average} & \multicolumn{2}{c|}{\textbf{45.7\%}} & \multicolumn{2}{c|}{\textbf{45.5\%}} & \multicolumn{2}{c|}{\textbf{73.0\%}} & \textbf{50.6\%} \\
\hline
\hline
\multicolumn{8}{|c|}{\textbf{Complete Pairing (BLS12\_381)}} \\
\hline
Single pair & 10,064 & 7,795 & 14,027 & 11,742 & 3,913 & 3,123 & 19.6\% \\
% n=2 pairs & 13,158 & 9,541 & 17,232 & 13,581 & 4,609 & 3,145 & 26.8\% \\
Three pairs & 16,484 & 11,287 & 20,669 & 15,420 & 5,421 & 3,167 & 32.7\% \\
\hline
\textbf{Average} & \multicolumn{2}{c|}{\textbf{31.5\%}} & \multicolumn{2}{c|}{\textbf{25.4\%}} & \multicolumn{2}{c|}{\textbf{41.6\%}} & \textbf{26.4\%} \\
\hline
\multicolumn{8}{l}{\footnotesize MULMOD/ADDMOD: Field operations in $\mathbb{F}_p$} \\
\multicolumn{8}{l}{\footnotesize POSEIDON: Hash operations for Fiat-Shamir; } \\
\multicolumn{8}{l}{\footnotesize Speedup: VM Cycle reduction.} \\
\hline
\end{tabular}
\end{table*}

\subsection{Comparison with State of the Art}

For reference, Gnark---currently the leading framework for field emulation in R1CS---reports 960,784 constraints for a single BN254 pairing and 311,672 for the Miller loop alone. Our Cairo implementation achieves 53,130 steps for a single Miller loop and 155,366 for a complete pairing, reflecting a fundamentally different algorithmic approach rather than just implementation tuning.

The key structural difference is that Gnark verifies each emulated $\Fe{12}$ multiplication individually inside the circuit, while our scheme batches all polynomial verification into a single bivariate check outside the inner circuit---with only the remainder polynomials committed to. This is why our commitment overhead (POSEIDON hashes) drops by 77\% while constraint counts drop by a more modest 45\%: we are eliminating a different class of cost.

Porting our IOP to Gnark is a direction for future work. The main challenge is that Gnark circuits allow only a single hash batching point, while our two-round IOP requires two sequential Fiat-Shamir challenges. One approach is an in-circuit GKR proof for the inner verification; we leave this as an open problem.

\subsection{Final Remarks}

We have shown that choosing the right granularity for polynomial identity testing---entire Miller loop steps rather than individual multiplications---yields substantial and measurable improvements. The key technique, collapsing a composite multi-operand product into a single Euclidean division check, is simple to implement and broadly applicable to other multi-step cryptographic computations in arithmetized environments.

Our implementation in Garaga is deployed on Starknet and has already enabled proof verification workloads previously infeasible on-chain. We hope this work encourages further exploration of polynomial batching granularity as a design axis for constraint-system implementations.


%%%% 8. BIBLIOGRAPHY %%%%
\bibliographystyle{alpha}
\bibliography{references}
%%%% NOTES
% - Consider using abbrev3.bib and crypto.bib from https://cryptobib.di.ens.fr/
% - Ensure all references are complete and properly formatted

\end{document}